\documentclass[12pt,a4paper]{article}
\usepackage[a4paper,left=2.5cm,right=2cm,top=2cm,bottom=2cm]{geometry}
\usepackage{fontspec}
\setmainfont{DejaVu Serif}
\setmonofont{DejaVu Sans Mono}
\usepackage{setspace}
\onehalfspacing
\usepackage{enumitem}
\usepackage{xcolor}
\usepackage{listings}
\usepackage{hyperref}
\hypersetup{hidelinks}
\usepackage{titlesec}
\titleformat{\section}{\normalfont\Large\bfseries}{\thesection.}{0.5em}{}
\titleformat{\subsection}{\normalfont\large\bfseries}{\thesubsection.}{0.5em}{}
\setlength{\parindent}{1.25cm}
\setlength{\parskip}{0pt}
\renewcommand{\contentsname}{Contents}
\definecolor{codebg}{RGB}{246,246,246}
\definecolor{codeframe}{RGB}{215,215,215}
\renewcommand{\lstlistingname}{Figure}
\lstset{basicstyle=\ttfamily\footnotesize,backgroundcolor=\color{codebg},frame=single,rulecolor=\color{codeframe},breaklines=true,columns=fullflexible,keepspaces=true,showstringspaces=false,aboveskip=0.7em,belowskip=0.7em,captionpos=b}
\begin{document}
\begin{titlepage}
\begin{center}
{\large Ministerul Educației și Cercetării al Republicii Moldova}\\[0.3cm]
{\large Universitatea Tehnică a Moldovei}\\[0.3cm]
{\large Facultatea Calculatoare, Informatică și Microelectronică}\\[0.3cm]
{\large Program: Software Engineering}\\[2.3cm]
{\Large \textbf{Operating Systems}}\\[0.8cm]
{\LARGE \textbf{Laboratory Work No. 9}}\\[0.5cm]
{\Large \textbf{Watching System Calls with strace}}\\[0.35cm]
{\large Track A: the real OS, from the outside (Linux)}
\vfill
\begin{flushright}
\begin{tabular}{ll}
\textbf{Elaborated by:} & \\
st. gr. FAF-241 & \\
\textbf{Ermicev Sofia} & \\
\\
\textbf{Verified by:} & \\
\textbf{Patricia Reițman} &
\end{tabular}
\end{flushright}
\vfill
{\large Chișinău -- 2026}
\end{center}
\end{titlepage}
\tableofcontents
\newpage
\section{Introduction}
The goal of this laboratory work was to observe the boundary between a Linux user-space program and the kernel using \texttt{strace}. I traced the system calls of an ordinary \texttt{ls -l} command, implemented a C program using \texttt{open()}, \texttt{write()}, \texttt{read()} and \texttt{close()}, and counted the resulting \texttt{openat()} calls.

All tests were performed on Ubuntu under WSL2 as user \texttt{sofia}, in \texttt{/home/sofia/os-lab9}. The installed \texttt{strace} version was 6.8. The terminal outputs reproduced below are from the actual executions performed for this laboratory; the complete traces were also saved as text files.

\begin{lstlisting}[caption={Environment check and strace version}]
$ whoami
sofia
$ pwd
/home/sofia/os-lab9
$ strace --version
strace -- version 6.8
\end{lstlisting}

\section{Part 1. See the Calls of a Real Command}
I ran \texttt{ls -l} under \texttt{strace}, restricting the displayed events to \texttt{openat}, \texttt{read}, \texttt{write}, and \texttt{close}. The \texttt{-f} option also requests tracing of any child processes.

\begin{lstlisting}[caption={Command used to inspect the system calls of ls}]
$ strace -f -e trace=openat,read,write,close ls -l
\end{lstlisting}

The trace included attempts to open shared libraries, locale files, \texttt{/proc/filesystems}, and the current directory. An \texttt{ENOENT} result on a locale lookup means that a candidate path was not found; this did not stop \texttt{ls} from running. At the time of the first run, the working directory was empty.

\begin{lstlisting}[caption={Selected lines from the real ls trace}]
openat(AT_FDCWD, "/etc/ld.so.cache", O_RDONLY|O_CLOEXEC) = 3
close(3)                                = 0
openat(AT_FDCWD, "/proc/filesystems", O_RDONLY|O_CLOEXEC) = 3
read(3, "nodev\tsysfs\nnodev\ttmpfs\nnodev\tbd"..., 1024) = 478
read(3, "", 1024)                       = 0
close(3)                                = 0
openat(AT_FDCWD, ".", O_RDONLY|O_NONBLOCK|O_CLOEXEC|O_DIRECTORY) = 3
close(3)                                = 0
write(1, "total 0\n", 8)                = 8
close(1)                                = 0
close(2)                                = 0
+++ exited with 0 +++
\end{lstlisting}

\subsection{Observe}
\textbf{Which system call does \texttt{ls} use to print to the screen? To which descriptor does it write?}

The output was printed using \texttt{write()}. The trace contains \texttt{write(1, "total 0\\n", 8) = 8}: descriptor \texttt{1} is standard output (\texttt{stdout}), and eight bytes were written successfully. Even a simple directory listing uses many system calls for initialization and accessing files, not only for printing the result.

\section{Part 2. A Program That Calls the Kernel Directly}
I created \texttt{syscalls.c} to open \texttt{note.txt}, write \texttt{hello, kernel} followed by a newline, close the file, reopen it read-only, and read the bytes back. The program uses the file-descriptor system-call interface rather than \texttt{fprintf} or \texttt{fread} for file contents.

\begin{lstlisting}[language=C,caption={Implementation of syscalls.c}]
#include <fcntl.h>
#include <unistd.h>
#include <string.h>
#include <stdio.h>

int main(void) {
    const char *msg = "hello, kernel\n";

    int fd = open("note.txt", O_CREAT | O_WRONLY | O_TRUNC, 0644);
    if (fd < 0) {
        perror("open");
        return 1;
    }

    printf("open returned descriptor %d\n", fd);

    ssize_t n = write(fd, msg, strlen(msg));
    printf("write sent %zd bytes\n", n);
    close(fd);

    fd = open("note.txt", O_RDONLY);
    char buf[64];
    n = read(fd, buf, sizeof(buf) - 1);
    buf[n] = '\0';

    printf("read got %zd bytes: %s", n, buf);
    close(fd);

    return 0;
}
\end{lstlisting}

\begin{lstlisting}[caption={Compilation and actual program output}]
$ gcc -Wall -Wextra -o syscalls syscalls.c
$ ./syscalls
open returned descriptor 3
write sent 14 bytes
read got 14 bytes: hello, kernel
$ cat note.txt
hello, kernel
\end{lstlisting}

The file contained the expected 14 bytes, including the final newline. I then traced the program's system calls.

\begin{lstlisting}[caption={Relevant events from the real syscalls trace}]
$ strace -e trace=openat,write,read,close ./syscalls
openat(AT_FDCWD, "note.txt", O_WRONLY|O_CREAT|O_TRUNC, 0644) = 3
write(3, "hello, kernel\n", 14)         = 14
close(3)                                = 0
openat(AT_FDCWD, "note.txt", O_RDONLY)  = 3
read(3, "hello, kernel\n", 63)          = 14
close(3)                                = 0
+++ exited with 0 +++
\end{lstlisting}

The trace also contained library initialization and \texttt{write(1, ...)} calls corresponding to the program's \texttt{printf()} messages. The displayed file operations show that Linux used \texttt{openat()} to implement the C \texttt{open()} requests.

\subsection{Observe}
\textbf{Why did the first \texttt{openat()} for \texttt{note.txt} return descriptor 3 rather than 0, 1, or 2?}

A process normally starts with the three standard descriptors already open: \texttt{0} for standard input (\texttt{stdin}), \texttt{1} for standard output (\texttt{stdout}), and \texttt{2} for standard error (\texttt{stderr}). Therefore the first free descriptor was \texttt{3}. Once it was closed, the same descriptor number could be reused for the second opening of \texttt{note.txt}.

\section{Part 3. Count One Kind of Call}
I restricted tracing to \texttt{openat}, saved the combined output, and used \texttt{grep} to count lines beginning with the system-call name. The \texttt{2>\&1} redirection matters because \texttt{strace} normally writes its trace to standard error.

\begin{lstlisting}[caption={Tracing and counting openat calls}]
$ strace -e trace=openat ./syscalls 2>&1 | tee openat_trace.txt
openat(AT_FDCWD, "/etc/ld.so.cache", O_RDONLY|O_CLOEXEC) = 3
openat(AT_FDCWD, "/lib/x86_64-linux-gnu/libc.so.6", O_RDONLY|O_CLOEXEC) = 3
openat(AT_FDCWD, "note.txt", O_WRONLY|O_CREAT|O_TRUNC, 0644) = 3
openat(AT_FDCWD, "note.txt", O_RDONLY)  = 3
open returned descriptor 3
write sent 14 bytes
read got 14 bytes: hello, kernel
+++ exited with 0 +++
$ grep -c '^openat(' openat_trace.txt
4
\end{lstlisting}

Of the four observed \texttt{openat()} calls, two were for initialization files (the dynamic linker cache and the C library), and two were for \texttt{note.txt}. This shows why the number of system calls made by a program can exceed the number explicitly written in its source code.

\section{Conclusion}
This laboratory helped me understand how normal C file operations appear at the kernel interface. I observed that \texttt{ls} writes to file descriptor \texttt{1}, while my own program uses descriptor \texttt{3} to access a file. The recorded trace confirmed the 14-byte write and read, and the final experiment counted four \texttt{openat()} calls, including startup activity that was not explicitly part of the application logic. The experiment also showed how \texttt{strace} can help inspect program behavior without changing the program itself.
\end{document}
